% ============================================================
%  CHUONG 2. NGHIEN CUU LY THUYET
% ============================================================
\chapter{Nghiên cứu lý thuyết}
\label{ch:ly-thuyet}

Chương này trình bày cơ sở lý thuyết của hai nhóm nội dung: nhóm công nghệ nền
tảng dùng để xây dựng hệ thống (mục 2.1 đến 2.4) và nhóm cơ sở thuật toán dùng để
giải các bài toán nghiệp vụ (mục 2.5). Với mỗi công nghệ, chương không dừng ở
việc mô tả lý thuyết mà còn có một mục riêng chỉ rõ công nghệ đó được vận dụng
vào đâu trong đồ án, để phần thiết kế ở Chương~\ref{ch:hien-thuc} có căn cứ tham
chiếu trực tiếp.

\section{Kiến trúc MVC và nền tảng ASP.NET Core}
\label{sec:mvc}

\subsection{Mẫu kiến trúc MVC}

Model - View - Controller (MVC) là mẫu kiến trúc phân chia ứng dụng thành ba
nhóm thành phần chính là Model, View và Controller, nhằm đạt được sự \textbf{phân
tách mối quan tâm} (separation of concerns) [12]. Theo mô tả của Microsoft, yêu
cầu từ người dùng được định tuyến tới Controller; Controller có nhiệm vụ làm việc
với Model để thực hiện hành động hoặc lấy kết quả truy vấn, sau đó chọn View để
hiển thị và cung cấp cho View dữ liệu cần thiết [12].

Vai trò của từng thành phần được xác định như sau [12]:

\begin{itemize}
  \item \textbf{Model} biểu diễn trạng thái của ứng dụng cùng các logic nghiệp vụ
        tác động lên trạng thái đó.
  \item \textbf{View} chịu trách nhiệm trình bày nội dung qua giao diện người
        dùng. Trong View chỉ nên chứa lượng logic tối thiểu và logic đó phải liên
        quan tới việc trình bày.
  \item \textbf{Controller} xử lý tương tác của người dùng, làm việc với Model và
        lựa chọn View để kết xuất.
\end{itemize}

Một đặc điểm quan trọng của mẫu này là \textbf{chiều phụ thuộc một chiều}: View
và Controller đều phụ thuộc vào Model, nhưng Model không phụ thuộc ngược lại vào
View hay Controller [12]. Nhờ đó Model có thể được xây dựng và kiểm thử độc lập
với phần trình bày, đây chính là tính chất mà đồ án khai thác để kiểm thử các
thuật toán lõi bằng kiểm thử đơn vị (trình bày ở Chương~\ref{ch:ket-qua}).

Tài liệu của Microsoft cũng khuyến nghị rằng Controller không nên gánh quá nhiều
trách nhiệm, và để tránh điều đó thì nên đẩy logic nghiệp vụ ra khỏi Controller
[12]. Khuyến nghị này là căn cứ trực tiếp cho quyết định thiết kế của đồ án: tách
riêng một tầng \textbf{Service} để chứa ba thuật toán lõi, còn Controller chỉ
đóng vai trò tiếp nhận yêu cầu, gọi Service và trả về View.

\subsection{Các thành phần của ASP.NET Core MVC được sử dụng}

ASP.NET Core MVC là khung làm việc mã nguồn mở, nhẹ và có tính kiểm thử cao dành
cho tầng trình bày [12]. Đồ án sử dụng trực tiếp các cơ chế sau của khung này:

\begin{itemize}
  \item \textbf{Định tuyến (routing):} ánh xạ URL tới hành động của Controller.
        Đồ án sử dụng định tuyến theo quy ước với khuôn dạng
        \code{\{controller\}/\{action\}/\{id?\}} [12].
  \item \textbf{Model binding:} tự động chuyển dữ liệu từ yêu cầu HTTP (giá trị
        biểu mẫu, dữ liệu tuyến, tham số chuỗi truy vấn) thành đối tượng mà
        Controller xử lý được [12].
  \item \textbf{Kiểm tra hợp lệ (model validation):} thực hiện thông qua các
        thuộc tính chú giải dữ liệu, được kiểm tra ở cả phía máy khách trước khi
        gửi và phía máy chủ trước khi hành động của Controller được gọi [12].
  \item \textbf{Tiêm phụ thuộc (dependency injection):} Controller yêu cầu các
        dịch vụ cần thiết thông qua hàm khởi tạo [12].
  \item \textbf{Bộ lọc (filters):} đóng gói các mối quan tâm cắt ngang như xử lý
        ngoại lệ hay phân quyền; thuộc tính \code{[Authorize]} chính là bộ lọc
        phân quyền có sẵn của khung [12].
  \item \textbf{Areas:} cơ chế phân hoạch ứng dụng lớn thành các nhóm chức năng
        nhỏ hơn [12].
  \item \textbf{Razor và Tag Helpers:} Razor là ngôn ngữ đánh dấu mẫu cho phép
        nhúng mã C\# vào HTML để sinh nội dung phía máy chủ; Tag Helpers cho phép
        mã phía máy chủ tham gia vào việc tạo và kết xuất phần tử HTML [12].
\end{itemize}

\subsection{Áp dụng trong đồ án}

Bảng~\ref{tab:mvc-ap-dung} đối chiếu từng cơ chế nêu trên với vị trí sử dụng cụ
thể trong hệ thống CookingAdvisor.

\begin{table}[H]
\centering
\caption{Các cơ chế của ASP.NET Core MVC và cách áp dụng trong đồ án}
\label{tab:mvc-ap-dung}
\begin{tabularx}{\textwidth}{p{3.4cm} Y}
\toprule
\bfseries Cơ chế & \bfseries Áp dụng cụ thể trong CookingAdvisor \\
\midrule
Định tuyến & Toàn bộ chức năng dùng định tuyến theo quy ước, ví dụ
  \code{/Recipe/Details/5}, \code{/Suggestion}, \code{/Menu/Generate} \\
\rowhl Model binding & Gom toàn bộ tham số tìm kiếm, lọc, sắp xếp và phân trang
  của trang danh sách món vào một lớp \code{RecipeFilterViewModel} duy nhất \\
Kiểm tra hợp lệ & Biểu mẫu đăng ký, đăng nhập và biểu mẫu thêm hoặc sửa món ăn
  trong khu quản trị \\
\rowhl Tiêm phụ thuộc & Năm lớp Service được đăng ký vòng đời \code{Scoped} và
  đưa vào Controller qua hàm khởi tạo \\
Bộ lọc & \code{[Authorize]} cho nhóm chức năng của Người dùng,
  \code{[Authorize(Roles = "Admin")]} cho toàn bộ khu quản trị \\
\rowhl Areas & Khu vực quản trị tách riêng thành \code{Areas/Admin} \\
Razor và Tag Helpers & Toàn bộ giao diện; ngoài ra đồ án tự viết Tag Helper
  riêng cho hệ thống biểu tượng, đặt trong thư mục \code{TagHelpers} \\
\bottomrule
\end{tabularx}
\end{table}

\section{Entity Framework Core và tiếp cận Code-First}
\label{sec:efcore}

\subsection{Vai trò của bộ ánh xạ đối tượng quan hệ}

Entity Framework Core (EF Core) là bộ ánh xạ đối tượng - quan hệ (Object
Relational Mapper) cho .NET, cho phép thao tác với cơ sở dữ liệu quan hệ thông
qua các đối tượng .NET thay vì viết câu lệnh SQL trực tiếp.

Ngoài lợi ích về năng suất, cách tiếp cận này còn mang lại một lợi ích về an toàn
thông tin thường bị bỏ qua: các truy vấn viết bằng LINQ được EF Core dịch thành
câu lệnh SQL có tham số hóa, nên giá trị do người dùng nhập không bao giờ được
ghép trực tiếp vào chuỗi lệnh. Nguy cơ chèn mã SQL vì vậy được loại bỏ ngay tại
tầng truy cập dữ liệu chứ không phụ thuộc vào kỷ luật của từng lập trình viên.

\subsection{Cơ chế migrations}

Đồ án sử dụng tiếp cận \textbf{Code-First}, trong đó lược đồ cơ sở dữ liệu được
sinh ra từ các lớp thực thể trong mã nguồn. Cơ chế bảo đảm sự đồng bộ giữa mã
nguồn và cơ sở dữ liệu là \textbf{migrations}. Tài liệu của Microsoft mô tả:
trong các dự án thực tế, mô hình dữ liệu thay đổi khi các tính năng được hiện
thực hóa, và lược đồ cơ sở dữ liệu cần thay đổi tương ứng để đồng bộ với ứng
dụng; tính năng migrations của EF Core cung cấp cách cập nhật lược đồ một cách
tăng dần nhằm giữ đồng bộ với mô hình dữ liệu trong khi vẫn bảo toàn dữ liệu hiện
có [11].

Cơ chế hoạt động của migrations gồm hai bước [11]:

\begin{enumerate}[label=\arabic*.]
  \item Khi mô hình thay đổi, lập trình viên dùng công cụ dòng lệnh của EF Core
        để tạo một migration mô tả các cập nhật cần thiết. EF Core so sánh mô
        hình hiện tại với \textbf{ảnh chụp} (snapshot) của mô hình cũ để xác định
        khác biệt và sinh ra các tệp mã nguồn migration; các tệp này được quản lý
        phiên bản như mọi tệp mã nguồn khác.
  \item Migration sau khi được sinh ra sẽ được áp dụng vào cơ sở dữ liệu. EF Core
        ghi nhận mọi migration đã áp dụng vào một \textbf{bảng lịch sử} đặc biệt,
        nhờ đó biết được migration nào đã và chưa được áp dụng.
\end{enumerate}

Hai lệnh chính được sử dụng là \code{dotnet ef migrations add <Tên>} để tạo
migration và \code{dotnet ef database update} để áp dụng vào cơ sở dữ liệu [11].

\subsection{Áp dụng trong đồ án}

Lựa chọn Code-First mang lại ba lợi ích thực tiễn được khai thác trực tiếp trong
đồ án.

Thứ nhất, hội đồng đánh giá có thể dựng lại toàn bộ cơ sở dữ liệu trên máy của
mình chỉ bằng một lệnh, không cần tệp sao lưu cơ sở dữ liệu và không phụ thuộc
vào phiên bản cụ thể của công cụ quản trị. Đây chính là nội dung của yêu cầu phi
chức năng về khả năng tái lập trình bày ở mục~\ref{sec:yeu-cau-phi-chuc-nang}.

Thứ hai, các quy tắc ràng buộc phức tạp như khóa chính ghép của bảng
\code{RecipeIngredient} và bảng \code{Favorite}, ràng buộc duy nhất trên cột
\code{MenuPlanId} của bảng \code{ShoppingList}, và hành vi xóa của từng quan hệ
đều được khai báo tập trung tại một chỗ là phương thức \code{OnModelCreating} của
lớp \code{AppDbContext}. Toàn bộ thiết kế này được trình bày chi tiết ở
mục~\ref{sec:mo-hinh-csdl}.

Thứ ba, do các tệp migration nằm trong hệ thống quản lý phiên bản cùng với mã
nguồn, lịch sử thay đổi lược đồ có thể truy vết được. Đồ án cũng xuất toàn bộ
chuỗi migration ra một script T-SQL đặt tại \code{setup/sql/InitialCreate.sql} để
phục vụ trường hợp máy đánh giá chưa cài công cụ dòng lệnh \code{dotnet-ef}.

\section{ASP.NET Core Identity}
\label{sec:identity}

\subsection{Các cơ chế Identity cung cấp}

ASP.NET Core Identity là hệ thống quản lý thành viên của nền tảng, cung cấp sẵn
các chức năng đăng ký, đăng nhập, lưu trữ thông tin người dùng và quản lý vai trò
[10]. Việc sử dụng Identity thay vì tự cài đặt cơ chế xác thực mang lại ba lợi
ích mà đồ án khai thác.

\textbf{Thứ nhất, về lưu trữ mật khẩu.} Identity không lưu mật khẩu dạng nguyên
bản mà lưu giá trị băm sinh bởi hàm băm mật khẩu chuyên dụng có kèm muối (salt).
Đây là yêu cầu bắt buộc về an toàn thông tin: nếu cơ sở dữ liệu bị lộ, kẻ tấn
công không thu được mật khẩu gốc của người dùng.

\textbf{Thứ hai, về chính sách mật khẩu.} Identity cho phép cấu hình các ràng
buộc như độ dài tối thiểu, yêu cầu chữ số, chữ hoa hay ký tự đặc biệt [10]. Đồ án
cấu hình độ dài tối thiểu là 8 ký tự.

\textbf{Thứ ba, về phân quyền theo vai trò.} Identity hỗ trợ mô hình vai trò
(role-based authorization), cho phép gán người dùng vào các vai trò và giới hạn
truy cập theo vai trò. Kết hợp với bộ lọc \code{[Authorize]} của MVC [12], hệ
thống có thể khai báo phân quyền ngay tại lớp Controller.

\subsection{Phân biệt xác thực và phân quyền}

Cần phân biệt rõ hai khái niệm thường bị nhầm lẫn: \textbf{xác thực}
(authentication) trả lời câu hỏi \textit{người dùng này là ai}, còn \textbf{phân
quyền} (authorization) trả lời câu hỏi \textit{người dùng này được phép làm gì}.
Một hệ thống chỉ kiểm tra xác thực mà bỏ qua phân quyền vẫn có lỗ hổng: người
dùng đã đăng nhập hợp lệ vẫn có thể truy cập tài nguyên không thuộc quyền của
mình.

Từ đó có thể phân biệt tiếp hai mức phân quyền. \textbf{Phân quyền mức chức năng}
kiểm tra người dùng có được phép gọi một hành động hay không, thường cài đặt bằng
thuộc tính khai báo trên Controller. \textbf{Phân quyền mức bản ghi} kiểm tra
người dùng có được phép thao tác trên \textit{đúng bản ghi đang xét} hay không;
mức này không thể giải quyết bằng thuộc tính khai báo mà phải đưa điều kiện lọc
vào chính câu truy vấn.

\subsection{Áp dụng trong đồ án}

Đồ án định nghĩa hai vai trò là \code{Admin} và \code{User}. Toàn bộ khu vực quản
trị được bảo vệ bằng \code{[Authorize(Roles = "Admin")]}, đó là phân quyền mức
chức năng.

Ở mức bản ghi, các thao tác trên thực đơn và danh sách đi chợ đều kiểm tra bản
ghi có thuộc về đúng người dùng đang đăng nhập hay không, và phép kiểm tra này
được đặt ngay trong biểu thức truy vấn chứ không đặt trong một câu lệnh điều kiện
sau khi đã nạp dữ liệu. Cách cài đặt cụ thể được phân tích ở
mục~\ref{sec:thuat-toan-di-cho}.

\section{Hệ quản trị cơ sở dữ liệu SQL Server và triển khai đa nền tảng}
\label{sec:sqlserver}

\subsection{Lý do lựa chọn SQL Server}

Đồ án sử dụng Microsoft SQL Server làm hệ quản trị cơ sở dữ liệu quan hệ, với các
lý do: khả năng tương thích tốt nhất với EF Core trong hệ sinh thái .NET, hỗ trợ
đầy đủ ràng buộc toàn vẹn tham chiếu cần thiết cho lược đồ nhiều bảng của đề tài,
và tính phổ biến trong môi trường đào tạo.

Yêu cầu về ràng buộc toàn vẹn tham chiếu không phải là yêu cầu hình thức. Như sẽ
trình bày ở mục~\ref{sec:hanh-vi-xoa}, đồ án cấu hình hành vi xóa khác nhau cho
từng quan hệ, trong đó một số quan hệ được đặt ở chế độ chặn xóa để bảo vệ dữ
liệu của người dùng khác. Cơ chế này chỉ có hiệu lực khi hệ quản trị thực thi
ràng buộc khóa ngoại ở mức máy chủ.

\subsection{Vấn đề đa nền tảng và giải pháp}

Một vấn đề thực tiễn phát sinh trong quá trình thực hiện là môi trường phát triển
và môi trường đánh giá không đồng nhất về hệ điều hành. Giải pháp được áp dụng
như sau:

\begin{itemize}
  \item \textbf{Trên macOS:} SQL Server được chạy trong \textbf{container} thông
        qua OrbStack hoặc Docker Desktop, sử dụng ảnh chính thức của Microsoft.
        Cấu hình container được mô tả trong tệp \code{docker-compose.yml} kèm
        theo mã nguồn.
  \item \textbf{Trên Windows:} SQL Server được cài đặt trực tiếp theo bản Express
        hoặc Developer.
\end{itemize}

Điểm cần nhấn mạnh là \textbf{sự khác biệt này không ảnh hưởng tới mã nguồn}. Ứng
dụng chỉ giao tiếp với cơ sở dữ liệu qua chuỗi kết nối; thay đổi giữa hai môi
trường chỉ là thay đổi giá trị chuỗi kết nối, không có bất kỳ đoạn mã nào phân
nhánh theo hệ điều hành. Mô hình triển khai tương ứng được thể hiện ở
Hình~\ref{fig:trien-khai}.

\subsection{Quản lý chuỗi kết nối}

Về an toàn thông tin, chuỗi kết nối chứa thông tin nhạy cảm nên \textbf{không
được lưu trong mã nguồn}. Đồ án sử dụng cơ chế User Secrets của .NET trong môi
trường phát triển; tệp \code{appsettings.json} được đưa vào hệ thống quản lý
phiên bản chỉ chứa giá trị rỗng làm chỗ giữ chỗ.

Cách làm này giữ cho kho mã nguồn công khai không chứa mật khẩu, đồng thời vẫn
cho phép mỗi máy tự khai báo chuỗi kết nối phù hợp với môi trường của mình mà
không cần sửa bất kỳ tệp nào đang được quản lý phiên bản.

\section{Cơ sở thuật toán}
\label{sec:co-so-thuat-toan}

\subsection{Độ đo tương đồng tập hợp cho bài toán gợi ý}
\label{sec:do-do}

Bài toán gợi ý theo nguyên liệu được phát biểu như sau. Gọi $A$ là tập nguyên
liệu mà người dùng đang có, $I_R$ là tập nguyên liệu cần thiết của món $R$. Cần
định nghĩa một hàm số đo mức độ phù hợp giữa $A$ và $I_R$ để làm cơ sở xếp hạng.

\textbf{Hệ số Jaccard.} Độ đo tương đồng tập hợp được sử dụng phổ biến nhất là hệ
số Jaccard, định nghĩa bằng tỉ số giữa lực lượng phần giao và lực lượng phần hợp
của hai tập hợp [8]:

\begin{equation}
J(A, I_R) = \frac{|A \cap I_R|}{|A \cup I_R|}
\label{eq:jaccard}
\end{equation}

\textbf{Lý do không sử dụng hệ số Jaccard trong đề tài.} Hệ số Jaccard là độ đo
\textit{đối xứng}, nghĩa là nó phạt cả phần thừa của $A$ lẫn phần thiếu của $I_R$.
Đặc tính này không phù hợp với ngữ nghĩa của bài toán. Xét ví dụ cụ thể: người
dùng khai báo đang có 40 nguyên liệu trong bếp, món $R$ chỉ cần 3 nguyên liệu và
cả 3 đều nằm trong số đó. Về mặt thực tế, đây là món \textbf{nấu được ngay}, cần
được xếp hạng cao nhất. Nhưng theo công thức Jaccard:

\begin{equation*}
J(A, I_R) = \frac{3}{40} = 0{,}075
\end{equation*}

giá trị này rất thấp và món ăn sẽ bị xếp hạng gần cuối. Nguyên nhân là Jaccard
coi 37 nguyên liệu dư của người dùng như một sự ``không tương đồng'', trong khi
trên thực tế việc có dư nguyên liệu hoàn toàn không phải là điểm trừ.

\textbf{Độ phủ.} Đề tài vì vậy sử dụng độ đo \textbf{không đối xứng} là độ phủ
(coverage), chỉ xét tỉ lệ nguyên liệu của món được đáp ứng:

\begin{equation}
\mathrm{coverage}(A, I_R) = \frac{|A \cap I_R|}{|I_R|}
\label{eq:coverage}
\end{equation}

Với ví dụ trên, $\mathrm{coverage} = 3/3 = 1{,}0$, phản ánh đúng thực tế là món
nấu được ngay. Độ phủ luôn nhận giá trị trong đoạn $[0, 1]$; giá trị bằng 1 tương
đương với điều kiện $I_R \subseteq A$, tức là tập nguyên liệu còn thiếu
$M = I_R \setminus A$ là tập rỗng. Bảng~\ref{tab:so-sanh-do-do} tóm tắt sự khác
biệt giữa hai độ đo trên chính ví dụ vừa nêu.

\begin{table}[H]
\centering
\caption{So sánh hệ số Jaccard và độ phủ trên cùng một tình huống}
\label{tab:so-sanh-do-do}
\begin{tabularx}{\textwidth}{p{3.6cm} Y Y}
\toprule
\bfseries Tiêu chí & \bfseries Hệ số Jaccard & \bfseries Độ phủ \\
\midrule
Tính đối xứng & Đối xứng, phạt cả phần dư của người dùng &
  Không đối xứng, chỉ xét phần nguyên liệu của món \\
\rowhl Mẫu số & $|A \cup I_R|$, phụ thuộc số nguyên liệu người dùng khai báo &
  $|I_R|$, chỉ phụ thuộc bản thân món ăn \\
Giá trị trong ví dụ ($|A| = 40$, $|I_R| = 3$, giao bằng 3) &
  $3/40 = 0{,}075$ & $3/3 = 1{,}0$ \\
\rowhl Thứ hạng thu được & Gần cuối danh sách & Cao nhất, đúng thực tế \\
Điều kiện nấu được ngay & Không biểu diễn trực tiếp được &
  Tương đương với giá trị bằng 1 \\
\bottomrule
\end{tabularx}
\end{table}

\textbf{Xếp hạng nhiều tiêu chí.} Chỉ dùng một mình độ phủ vẫn chưa đủ, vì hai
món có cùng độ phủ nhưng khác nhau về số nguyên liệu còn thiếu sẽ không phân biệt
được. Đề tài sử dụng thứ tự từ điển trên bộ bốn tiêu chí, xét lần lượt:

\begin{enumerate}[label=\arabic*.]
  \item Món nấu được ngay ($M = \emptyset$) xếp trước.
  \item Độ phủ giảm dần.
  \item Số nguyên liệu còn thiếu $|M|$ tăng dần.
  \item Tên món theo thứ tự bảng chữ cái.
\end{enumerate}

Tiêu chí thứ tư có vai trò kỹ thuật quan trọng: nó bảo đảm \textbf{thứ tự tất
định} (deterministic). Nếu thiếu tiêu chí phá vỡ thế cân bằng này, hai lần chạy
cùng một truy vấn có thể cho ra thứ tự khác nhau, gây khó khăn cho việc kiểm thử
tự động và tạo trải nghiệm không nhất quán cho người dùng.

\textbf{Độ phức tạp.} Gọi $n$ là số món trong kho dữ liệu và $k$ là số nguyên
liệu trung bình của một món. Nếu biểu diễn $A$ bằng bảng băm, việc kiểm tra một
nguyên liệu có thuộc $A$ hay không tốn thời gian trung bình $O(1)$. Khi đó chi
phí tính toán độ phủ cho toàn bộ kho dữ liệu là $O(nk)$, cộng thêm $O(n \log n)$
cho bước sắp xếp, tổng cộng là $O(nk + n \log n)$.

\subsection{Thuật toán tham lam có ràng buộc cho bài toán lập thực đơn}
\label{sec:tham-lam}

\textbf{Phát biểu bài toán.} Cần gán món ăn vào $D \times M$ ô, với $D = 7$ ngày
và $M = 3$ bữa, tức 21 suất ăn, thỏa mãn các ràng buộc:

\begin{itemize}
  \item \textbf{Ràng buộc cứng:} món được gán phải thuộc vùng miền người dùng
        chọn (nếu có).
  \item \textbf{Ràng buộc mềm 1:} món phải phù hợp với bữa được gán (sáng, trưa
        hoặc tối).
  \item \textbf{Ràng buộc mềm 2:} hạn chế lặp món trong cùng một tuần.
  \item \textbf{Tiêu chí tối ưu:} tổng năng lượng trong ngày tiệm cận mức mục
        tiêu.
\end{itemize}

Bài toán tối ưu toàn cục với đầy đủ các ràng buộc trên thuộc lớp bài toán khó.
Tuy nhiên, với quy mô thực tế của đề tài (21 suất, vài chục món), lời giải tối ưu
tuyệt đối không phải là yêu cầu bắt buộc: người dùng chỉ cần một thực đơn hợp lý
và luôn có thể chỉnh sửa thủ công.

\textbf{Lựa chọn chiến lược tham lam.} Thuật toán tham lam là chiến lược tại mỗi
bước chọn phương án có vẻ tốt nhất ở thời điểm hiện tại, không quay lui để xét
lại các lựa chọn đã thực hiện [6]. Chiến lược này không bảo đảm nghiệm tối ưu
toàn cục trong trường hợp tổng quát, nhưng cho lời giải chấp nhận được với chi
phí tính toán thấp, phù hợp với yêu cầu phản hồi tức thời của ứng dụng web.

\textbf{Thứ tự nới lỏng ràng buộc.} Điểm thiết kế quan trọng nhất của thuật toán
là quy định \textbf{thứ tự nới lỏng} khi tập ứng viên trở nên rỗng. Thứ tự này
được xác định theo mức độ ảnh hưởng tới trải nghiệm người dùng, từ nhẹ tới nặng:

\begin{enumerate}[label=\arabic*.]
  \item Ban đầu, tập ứng viên gồm các món hợp bữa và \textbf{chưa được dùng}
        trong tuần.
  \item Nếu tập này rỗng, nới ràng buộc \textbf{không lặp}, cho phép dùng lại món
        đã dùng nhưng vẫn giữ điều kiện hợp bữa.
  \item Nếu vẫn rỗng, nới tiếp ràng buộc \textbf{hợp bữa}.
  \item Ràng buộc \textbf{vùng miền không bao giờ được nới}, vì đó là lựa chọn
        tường minh của người dùng: một thực đơn được yêu cầu là món miền Nam mà
        lại chứa món miền Bắc sẽ bị xem là sai chức năng, chứ không phải là một
        sự linh hoạt.
\end{enumerate}

\textbf{Vấn đề suy thoái tính đa dạng.} Trong quá trình kiểm thử, một khiếm
khuyết đã được phát hiện ở bước 2. Khi buộc phải lặp món, nếu tiêu chí lựa chọn
vẫn là ``gần mức năng lượng mục tiêu nhất'', thì do trạng thái năng lượng tích
lũy trong ngày lặp lại theo chu kỳ, thuật toán có xu hướng chọn \textbf{đúng một
món} cho mọi suất còn trống. Kết quả đo được trên dữ liệu thử nghiệm là một món
xuất hiện 12 lần trong khi nhiều món khác không được sử dụng lần nào.

\textbf{Giải pháp.} Bổ sung một tiêu chí ưu tiên đứng trước tiêu chí năng lượng,
chỉ áp dụng trong nhánh phải lặp: \textbf{ưu tiên món có số lần đã sử dụng ít
nhất}. Về bản chất, đây là chuyển từ chiến lược tham lam thuần túy sang chiến
lược cân bằng tải theo tần suất sử dụng. Sau khi áp dụng, thực đơn sinh ra đạt 17
món khác nhau trên 21 suất. Chi tiết kiểm chứng được trình bày ở
Chương~\ref{ch:ket-qua}.

\textbf{Độ phức tạp.} Với $S = D \times M = 21$ suất và $n$ món ứng viên, mỗi
suất cần duyệt và sắp xếp tập ứng viên, cho độ phức tạp $O(S \cdot n \log n)$.
Với quy mô thực tế, chi phí này là không đáng kể.

\subsection{Bài toán gộp nhóm cho danh sách đi chợ}
\label{sec:gop-nhom}

Bài toán sinh danh sách đi chợ là bài toán \textbf{gộp nhóm và tổng hợp}
(grouping and aggregation) quen thuộc trong xử lý dữ liệu. Cho thực đơn $P$ gồm
các suất ăn, mỗi suất tham chiếu tới một món, mỗi món có danh sách bộ ba (nguyên
liệu, khối lượng, đơn vị). Cần tính:

\begin{equation}
Q(i, u) = \sum_{\substack{s \in P \\ (i, q, u) \in R_s}} q
\label{eq:gop-nhom}
\end{equation}

trong đó $R_s$ là tập nguyên liệu của món ở suất $s$.

Có hai điểm cần lưu ý về mặt thiết kế.

\textbf{Thứ nhất, khóa gộp nhóm phải là cặp (nguyên liệu, đơn vị) chứ không chỉ
là nguyên liệu.} Cùng một nguyên liệu có thể được định lượng bằng các đơn vị khác
nhau ở những công thức khác nhau, ví dụ hành lá tính theo gam ở món này nhưng
tính theo bó ở món khác. Cộng dồn hai giá trị khác đơn vị sẽ cho kết quả vô
nghĩa. Việc đưa đơn vị vào khóa gộp nhóm bảo đảm chỉ những giá trị cùng đơn vị
mới được cộng với nhau.

\textbf{Thứ hai, phép duyệt phải theo suất ăn chứ không theo món.} Nếu một món
xuất hiện ba lần trong tuần thì nguyên liệu của nó phải được tính ba lần. Nói
cách khác, phép duyệt thực hiện trên tập các suất ăn, không phải trên tập các món
khác nhau.

Độ phức tạp của thuật toán là $O(S \cdot k)$ với $S$ là số suất ăn và $k$ là số
nguyên liệu trung bình mỗi món, khi sử dụng bảng băm cho thao tác gộp nhóm.

\subsection{Tổng hợp cơ sở thuật toán của ba chức năng lõi}

Bảng~\ref{tab:tong-hop-thuat-toan} tổng hợp ba thuật toán vừa trình bày, làm căn
cứ cho phần cài đặt ở mục~\ref{sec:cai-dat-thuat-toan}.

\begin{table}[H]
\centering
\caption{Tổng hợp cơ sở thuật toán của ba chức năng lõi}
\label{tab:tong-hop-thuat-toan}
\begin{tabularx}{\textwidth}{p{2.9cm} Y Y p{2.9cm}}
\toprule
\bfseries Chức năng & \bfseries Lớp bài toán & \bfseries Kỹ thuật sử dụng &
\bfseries Độ phức tạp \\
\midrule
Gợi ý theo nguyên liệu & Xếp hạng theo độ tương đồng tập hợp &
  Độ phủ không đối xứng, sắp xếp theo thứ tự từ điển bốn tiêu chí &
  $O(nk + n \log n)$ \\
\rowhl Sinh thực đơn tuần & Xếp lịch có ràng buộc &
  Tham lam có thứ tự nới lỏng xác định, cân bằng theo tần suất sử dụng &
  $O(S \cdot n \log n)$ \\
Sinh danh sách đi chợ & Gộp nhóm và tổng hợp &
  Bảng băm với khóa ghép (nguyên liệu, đơn vị) & $O(S \cdot k)$ \\
\bottomrule
\end{tabularx}
\end{table}

\section{Kết luận chương}

Chương 2 đã trình bày hai nhóm cơ sở lý thuyết của đồ án.

Về công nghệ, mẫu kiến trúc MVC cùng nguyên tắc phân tách mối quan tâm là căn cứ
cho quyết định tách tầng Service khỏi Controller, nhờ đó các thuật toán lõi có
thể được kiểm thử độc lập với giao diện. EF Core theo tiếp cận Code-First cùng cơ
chế migrations bảo đảm lược đồ cơ sở dữ liệu luôn đồng bộ với mã nguồn và có thể
tái tạo trên máy khác chỉ bằng một lệnh. ASP.NET Core Identity cung cấp nền tảng
xác thực và phân quyền theo vai trò đã được kiểm chứng, tránh được các sai sót
thường gặp khi tự cài đặt.

Về thuật toán, chương đã lập luận cho ba lựa chọn thiết kế then chốt: dùng độ phủ
thay cho hệ số Jaccard do tính chất không đối xứng phù hợp với ngữ nghĩa bài
toán; dùng chiến lược tham lam có thứ tự nới lỏng ràng buộc xác định cho bài toán
lập thực đơn, kèm cơ chế cân bằng tần suất để bảo toàn tính đa dạng; và dùng khóa
gộp nhóm ghép cặp (nguyên liệu, đơn vị) để bảo đảm tính đúng đắn của phép cộng
dồn khối lượng.

Các cơ sở lý thuyết này sẽ được vận dụng trực tiếp vào phần thiết kế và cài đặt ở
Chương~\ref{ch:hien-thuc}.
